\documentclass[11pt]{article}
\usepackage[margin=0.9in]{geometry}
\usepackage{booktabs,longtable,amsmath,amssymb,hyperref}
\usepackage[T1]{fontenc}
\hypersetup{colorlinks=true,linkcolor=blue,urlcolor=blue}
\title{\textbf{Disposition of the Consolidated Manuscript Error\\and Revision Register}}
\author{Karthikeya \\ \small Accretion Labs \\ \small \texttt{founder@accretionlabs.net}}
\date{11 October 2026}
\begin{document}
\maketitle

\section*{What this document is}

The register of 11 October 2026 lists 52 items. This document records the disposition of
each, and sets out the reasoning for every item that was \emph{not} applied as written. It
covers Draft~17 of \emph{Chemistry-resolved disagreement between PBE and r$^2$SCAN
stability verdicts in an AI-generated materials database}.

Of the 52 items, 45 are applied in full, 6 are applied in part, and 1 cannot be done by the
author alone. Nothing is declined on the ground that the register is wrong about the facts.
Where an item is not applied as written, it is for one of three reasons: the register itself
offered an alternative and the alternative was taken; the work the item asks for would churn
every table in the paper to change a number in the fourth significant figure; or the item
asks for something only a second person can supply.

Six items were resolved by measurement rather than by rewording, and in two cases the
measurement contradicted what the paper previously said. Those are listed in
Section~\ref{sec:measured} because they are the places where the register changed a result
and not only a sentence.

\section{Items not applied as written}

\subsection*{Item 01 --- the register's second option was taken, not its first}

The register offers two routes. The first regenerates every full-release analytical output
at $n = 139{,}516$; the second withdraws the universal-exclusion claim and labels the
retained outcomes correctly.

Route two was taken, and the reason is arithmetic. The 25 matched entries that were not
PBE-stable were counted: 24 of them fall inside a primary composition cell. Excluding them
moves the release-based rate from 13.0413\% to 13.0401\% and the primary-cell rate from
12.9150\% to 12.9136\%. Regenerating 40 cell rates, 83 element rates, every cluster
bootstrap interval, the regression, the band-gap table and the budget evaluation to move the
third decimal place of two percentages would create a large diff with no substantive change,
and would break the correspondence between the paper's $n$ and the released file that
readers will download.

Instead Section~3.4 now defines two outcomes separately --- \emph{release-based
disagreement}, which makes no use of the PBE verdict and is defined for all 139{,}541
entries, and a \emph{confirmed transition}, which requires an observed non-positive PBE
decomposition energy and is defined for the 117{,}018 matched PBE-stable entries. The 25 are
excluded wherever a confirmed transition is the outcome or an observed PBE energy is
required (the budget evaluation, the propensity-weighted rates, the energy-displacement
analysis, the matched comparison) and retained in the release-based rates and the cell
tabulations. Both rates are printed, with and without them. The universal-exclusion claim is
gone. The population-flow table now marks which rows are raw release counts and which are
analytical, and \texttt{run\_all.py} checks the 25 and the 24 as headline quantities so the
distinction cannot silently rot.

\subsection*{Item 10 --- seed stability was demonstrated instead of raising the replicate count}

The register permits either. Raising 300--600 replicates to, say, 5{,}000 would multiply the
runtime of the whole bootstrap stack, and would not answer the question the item actually
raises, which is whether the \emph{thresholded count} of significant contrasts is stable. So
the count was recomputed under five seeds: 53, 53, 52, 53, 53. That is reported in
Appendix~C and checked in CI as two headline quantities (minimum 52, maximum 53). If the
range had been wider the replicate count would have had to go up; it is not.

\subsection*{Item 11 --- the reduced-stoichiometry sensitivity check was not run}

The register marks it optional. The limitation is now stated explicitly in the paper:
compositions are used exactly as released and are not reduced, so the clustering key depends
on the released unit-cell convention as well as on the underlying stoichiometry, and the
same compound released as A$_2$B$_4$ and as AB$_2$ would receive different anonymised
patterns. The space-group source is named and the absence of missing values recorded.

Running the check would mean re-deriving all 8{,}429 cluster keys and re-running every
cluster bootstrap, cluster-robust regression and cluster-wise budget split on the new keys.
That is a day of compute and a second set of every interval in the paper, to test a proxy
that the paper already declines to claim a direction for. It is recorded as the obvious next
robustness check rather than done.

\subsection*{Item 47 --- independent review of the core definitions and scripts is not something I can supply}

Everything mechanical in this item is done: the pipeline was regenerated on final inputs,
script-to-output mappings were checked, the CI checks were extended from 23 headline
quantities to 38 (including the updated Benjamini--Hochberg count, the analytical population
counts and the new budget contrasts), the environment is pinned, and SHA-256 checksums of
both input files are now printed in Appendix~F.

The last clause --- \emph{obtain independent review of core definitions and scripts} ---
cannot be satisfied from inside the project. The Data and code availability section now says
so in as many words: the core definitions and scripts have not been reviewed by anyone other
than the author. This is the one item on the register that is open, and it is the one the
accompanying strategic note is right to call the highest-value remaining step.

\subsection*{Item 51 --- the citation ledger is partial, and the disclosure now says exactly how partial}

Six sources were checked against their source text rather than against the citation, and the
checks are recorded in the Disclosures section: Kingsbury et al.\ (6{,}307 materials behind
``approximately 6000 solids''; r$^2$SCAN more accurate than both SCAN and PBEsol for
formation energies); Kothakonda et al.\ (GGAs comparable for transition metals,
intermetallics and weakly bound solids; PBE slightly more accurate for the 233 Type~2
decomposition enthalpies); An et al.\ (603 to 298, 166 dynamically stable, 109 surviving,
SCAN and not r$^2$SCAN, and the exclusion of 142 rare-earth compounds for $f$-electron and
pseudopotential reasons); Yu et al.\ (19{,}998 unary and binary materials, per-element
formation-energy error); Huang et al.\ (380{,}190 paired MatPES entries, and thermodynamic
validation on exactly six binary systems: Fe--O, Mn--O, Ti--Cl, Mo--Cl, V--S, Cu--S); and
Liu and Meng, whose ``more than 60\% metallic'' figure is now cited as reported and used in
no analysis.

The remaining references were verified for existence, authorship and venue but not claim by
claim, and updated publication metadata was not re-checked for all of them. The paper does
not claim otherwise. Completing the ledger is a reading task on roughly 25 papers; it is
listed in Section~\ref{sec:yours}.

\subsection*{Item 52 --- the abstract was shortened; the supplement migration and the contents page were not done}

The abstract now fits on one page. Repeated defensive phrasing was cut in several places,
and the rhetorical formulations the register names (``not additive bookkeeping'', ``does not
generate'', ``strongest predictor in every specification'') are gone from the abstract,
results and conclusion.

Two sub-requests were not acted on. The design-effect arithmetic stays in the main text,
because the whole inferential apparatus of the paper rests on the claim that these
observations are not independent, and a reader who cannot follow that arithmetic in place
has no reason to accept any interval in the paper. The table of contents stays, because the
paper now carries eight appendices and a reader looking for the input manifest or the
all-cells table needs a way to find them.

The paper is also longer than the 39 pages the register saw, not shorter: it is 46. The
additions are the weight diagnostics the register demands in item 32, the all-cells
supplement it demands in item 38, the input manifest it demands in item 06, the model
specification it demands in item 16, and the stronger budget comparison. Those cost about
six pages between them. Length was traded for the evidence the register asked for; the
trade seems like the right one, but it is a trade and not a clean win.

\section{Items the register asked for that only you can do}\label{sec:yours}

None of these is a judgement call. They need an account, a person, or an email that is not
mine to send.

\begin{enumerate}
\item \textbf{Independent computational review} (item 47, and \S2 of the strategic note).
Give a reviewer the repository, the checksums, the pinned environment and the README, and
ask for two separate things: whether the central tables regenerate on a clean machine, and
whether the eligibility definitions, group encodings, bootstrap contrasts and budget splits
are \emph{correct}. The second is not the same as the first and is the one that matters.
\item \textbf{Ask the GNoME authors what the ``All'' column is} (item 41, \S3 of the
strategic note). The paper's reading of it is marked as an inference throughout and every
conclusion depending on it is flagged conditional, which is the best that can be done
without an answer. One email settles it.
\item \textbf{Deposit the repository in a DOI-minting archive} (Zenodo or similar) and cite
the DOI. The paper currently says the deposit has not been made.
\item \textbf{Complete the citation ledger} (item 51) on the roughly 25 references not yet
checked claim by claim.
\item \textbf{Regenerate the Materials Project API key.} It was partly visible in a
screenshot during this work.
\item \textbf{Make the repository public} before submission. The paper gives its URL.
\end{enumerate}

\section{Where the register changed a result, not a sentence}\label{sec:measured}

Six items were answered by running something. Two of them contradicted the paper.

\paragraph{Item 18 falsified a caption.} The caption of the logistic-regression table
claimed that dropping hydrides ``moves no coefficient in the third decimal''. Measured, it
moves the $U$-list odds ratio from 3.448 to 3.468 --- the third decimal of the largest term
in the table. The caption now prints all three measured shifts. The claim was wrong and had
never been computed.

\paragraph{Item 32 found that a trimming choice was load-bearing and undocumented.} Appendix
C said the inverse-probability weights were ``untrimmed''. The code floors the fitted
propensity at $10^{-2}$. That is not a cosmetic difference: 26 eligible recomputed entries
fall below the floor, and leaving them untrimmed drives the largest weight to
$6.6\times10^{5}$, collapses Kish's effective sample size from 83{,}875 to 5, and moves the
weighted rate from 12.71\% to 3.05\%. The reported 12.7\% is conditional on trimming, and
the paper now says so, reports the full diagnostics (propensity quantiles, overlap between
recomputed and unrecomputed, maximum weight, effective sample size, share of weight in the
heaviest 1\%) and shows the estimate is stable at 12.71\%, 12.74\% and 12.76\% for floors of
$10^{-2}$, $2\times10^{-2}$ and $5\times10^{-2}$.

\paragraph{Item 37 replaced a hand-wave with a standardisation.} The aggregate inversion was
previously explained by a commentary's ``more than 60\% metallic'' figure. The anion-class
mix is now measured directly: the release is 38.9\% intermetallic and 9.8\% oxide. Over the
twenty matched cells, standardising both sides to the known-materials cell mix, GNoME
disagrees \emph{more} by $+1.2$ points $[+0.5, +1.9]$; at the GNoME mix, by $+1.8$. The
direction of the aggregate comparison therefore reverses once the mix is held common, which
identifies composition as a sufficient explanation of the inversion over those cells.

\paragraph{Item 40 was answered by deletion, with evidence.} The claim that known-materials
cells concentrate within 20 meV/atom of the hull is gone. The Materials Project
\texttt{energy\_above\_hull} field is floored at zero and reads zero for 21.8\% of the
152{,}426 entries pulled, so those margins are not recoverable from it. The paper now says
that, and keeps the GNoME-side near-hull gradient as an internal descriptive result only.

\paragraph{Item 43 corrected a cluster count.} The main text said half of 8{,}429 clusters.
The code applies eligibility first; the eligible pool is 7{,}822 clusters and 3{,}911 are
drawn per split. Both numbers are now in the main text and both are checked in CI.

\paragraph{Item 44 was tested rather than hedged.} The register asked for the upper-estimate
language to be removed because the direction of generated-regressor bias is not established.
It is now removed, and the question was also settled empirically. On 20 fresh cluster-wise
splits, a hull-distance model with twenty quantile-bin dummies gains nothing over simple
closest-first ranking ($-0.06$ points at a 20\% budget, ahead in 3 splits of 20); the
combined policy beats that flexible model by $+6.00$ points (20 of 20); and cross-fitting
the cell rates leaves the margin at $+6.31$ points (20 of 20). Composition still adds after
hull distance is modelled nonparametrically, and the gain is not generated-regressor
optimism.

\section{Verification state of Draft 17}

Two layers, both green:

\begin{itemize}
\item \texttt{run\_all.py} regenerates the results directory from the two input CSVs and
checks 38 headline quantities, up from 23. Four are claim-level rather than value-level: the
budget ordering holds at every budget; the combined policy's margin over the flexible hull
model stays positive under cross-fitting; the iron excess is smallest in oxides of all anion
classes; composition adds after hull distance is modelled nonparametrically.
\item \texttt{scripts/check\_paper.py} parses every number out of the \LaTeX{} tables and
diffs 232 printed values against the result files.
\end{itemize}

Six new analysis scripts (\texttt{phase23} through \texttt{phase28}) are wired into the
driver and the continuous-integration workflow, so the register's diagnostics are regenerated
on every push rather than recorded once. The workflow deletes the stored results first,
exempting only the two cached Materials Project pulls, which require a credential.

The document compiles with no errors, no overfull or underfull boxes and no undefined
references, at 46 pages.

\section{Full disposition}

\begin{longtable}{rp{0.56\textwidth}l}
\toprule
Item & Register heading & Disposition \\
\midrule
\endfirsthead
\toprule
Item & Register heading & Disposition \\
\midrule
\endhead
01 & Population accounting consistency & Applied, second option \\
02 & Unmatched entries not bounded by 25 & Applied \\
03 & Release-based vs confirmed transitions & Applied \\
04 & Gap and cell counts under exclusion & Applied, verified \\
05 & Eligible sample for the IPW rate & Applied, verified \\
06 & Input manifest & Applied, with checksums \\
07 & Sandwich notation & Applied \\
08 & Joint bootstrap, 51 to 53 & Applied, verified \\
09 & Cluster pool per analysis & Applied \\
10 & Replicate or seed stability & Applied, seed stability \\
11 & Unreduced compositions & Applied; check not run \\
12 & Extremum selection is descriptive & Applied \\
13 & Exact boolean group definitions & Applied \\
14 & Pseudo-$R^2$ is not variance explained & Applied \\
15 & Marginal effects: definition and uncertainty & Applied, intervals computed \\
16 & Atomic-fraction interaction model & Applied \\
17 & Additivity constrains, does not require & Applied, heading changed \\
18 & Hydride caption precision claim & Applied, measured \\
19 & $U$ is not a sufficient explanation & Applied \\
20 & Causal verbs & Applied \\
21 & $U$ exclusivity claim & Applied \\
22 & Open-shell wording and copper & Applied \\
23 & Reference convention vs gauge shift & Applied \\
24 & Mean and dispersion both differ & Applied \\
25 & Frozen-$4f$ as a conditional concern & Applied \\
26 & Accuracy does not propagate into disagreement & Applied \\
27 & Shared selection remains unresolved & Applied \\
28 & Winner's curse is a mechanism, not a fact & Applied \\
29 & No mediation or independence claims & Applied \\
30 & Exact-zero gap provenance & Applied, verified at source \\
31 & IPW is a sensitivity check & Applied \\
32 & Weight and positivity diagnostics & Applied, measured \\
33 & Crystal-system adjustment claim & Applied \\
34 & Pseudopotential exclusion claim & Applied \\
35 & Per-feature attribution of weighting & Applied \\
36 & Entry coverage is a proxy & Applied \\
37 & Class proportions and standardisation & Applied, measured \\
38 & All matched cells in a supplement & Applied, new appendix \\
39 & One-sided bound needs its condition & Applied \\
40 & Known-material margins within 20 meV & Applied, deleted with evidence \\
41 & Unidentified alternative reference set & Applied \\
42 & Threshold definition and pnictide values & Applied \\
43 & Cluster pool for the budget splits & Applied, corrected \\
44 & Upper-estimate language & Applied, and cross-fitted \\
45 & Informal heuristics not evaluated & Applied \\
46 & Catch definition and split variability & Applied \\
47 & Clean regeneration and independent review & Applied in part; review open \\
48 & ``Strongest in every specification'' & Applied \\
49 & Coverage, metastability, Q1/Q3 framing & Applied \\
50 & Intercept label, plot reference, contrast & Applied \\
51 & Claim-level citation verification & Applied in part; 6 of $\sim$31 \\
52 & Abstract and front-matter length & Applied in part \\
\bottomrule
\end{longtable}

\end{document}
